% ==============================================================================
% CONFIGURACIÓN Y ESTILO VISUAL — METODOLOGÍAS DE DESARROLLO ÁGIL (MDA)
% Departamento de Lenguajes y Sistemas Informáticos — Universidad de Granada
% Plataforma repasaYA
% ==============================================================================

\usepackage[utf8]{inputenc}
\usepackage[T1]{fontenc}
\usepackage[spanish,es-tabla]{babel}
\usepackage{lmodern}
\usepackage{microtype}
\usepackage[dvipsnames,svgnames,table,xcdraw]{xcolor}
\usepackage{amsmath,amssymb,amsfonts,amsthm}
\usepackage{graphicx}
\usepackage{float}
\usepackage{booktabs}
\usepackage{tabularx}
\usepackage{longtable}
\usepackage{array}
\usepackage{multirow}
\usepackage{multicol}
\usepackage{eurosym}
\usepackage{enumitem}
\usepackage{fancyhdr}
\usepackage{titlesec}
\usepackage{listings}
\usepackage{tcolorbox}
\tcbuselibrary{skins,breakable}
\usepackage{tikz}
\usetikzlibrary{babel,arrows.meta,positioning,calc,shapes.geometric,fit,backgrounds}

% Configuración de márgenes
\usepackage[a4paper,top=2.5cm,bottom=2.5cm,left=2.6cm,right=2.6cm,headheight=24pt]{geometry}

% Paleta de colores profesionales para Ingeniería del Software / Ágil
\definecolor{mdaDarkBlue}{RGB}{18, 32, 58}
\definecolor{mdaPrimary}{RGB}{30, 74, 138}
\definecolor{mdaAccent}{RGB}{99, 91, 255}
\definecolor{mdaLightAccent}{RGB}{242, 243, 255}
\definecolor{mdaGreen}{RGB}{22, 135, 70}
\definecolor{mdaLightGreen}{RGB}{240, 253, 244}
\definecolor{mdaOrange}{RGB}{217, 119, 6}
\definecolor{mdaLightOrange}{RGB}{255, 251, 235}
\definecolor{mdaRed}{RGB}{185, 28, 28}
\definecolor{mdaLightRed}{RGB}{254, 242, 242}
\definecolor{mdaGray}{RGB}{246, 248, 251}
\definecolor{mdaBorderGray}{RGB}{215, 222, 232}
\definecolor{mdaTeal}{RGB}{15, 118, 110}
\definecolor{mdaYellowCard}{RGB}{255, 252, 220}
\definecolor{mdaYellowBorder}{RGB}{210, 185, 80}
\definecolor{codeBg}{RGB}{248, 249, 252}

% Hipervínculos
\usepackage{hyperref}
\hypersetup{
    colorlinks=true,
    linkcolor=mdaDarkBlue,
    citecolor=mdaTeal,
    urlcolor=mdaAccent,
    pdfauthor={repasaYA - Universidad de Granada (UGR)},
    pdftitle={Metodologías de Desarrollo Ágil - Seminario 2: Historias de Usuario},
    pdfsubject={Desarrollo Ágil de Software, Scrum, XP, Historias de Usuario, INVEST, Planning Poker, Spikes},
    pdfkeywords={MDA, Historias de Usuario, Scrum, XP, INVEST, CCC, BDD, Planning Poker, Velocidad, Spikes, UGR}
}

% Configuración de Encabezados y Pies de página
\pagestyle{fancy}
\fancyhf{}
\fancyhead[LE,RO]{\footnotesize\textsf{\thepage}}
\fancyhead[LO]{\footnotesize\textsf{\rightmark}}
\fancyhead[RE]{\footnotesize\textsf{\leftmark}}
\fancyfoot[C]{\footnotesize\textsf{\color{gray}repasaYA $\cdot$ Metodologías de Desarrollo Ágil (UGR) $\cdot$ Seminario 2: Historias de Usuario}}
\renewcommand{\headrulewidth}{0.4pt}
\renewcommand{\footrulewidth}{0.2pt}

% Configuración de títulos de capítulos y secciones
\titleformat{\chapter}[display]
  {\normalfont\bfseries\color{mdaDarkBlue}\sffamily}
  {\Large\textsf{Bloque \thechapter}}{1ex}
  {\titlerule[1.5pt]\vspace{1ex}\huge}
  [\vspace{1ex}\titlerule]

\titleformat{\section}
  {\normalfont\Large\bfseries\color{mdaPrimary}\sffamily}
  {\thesection}{0.8em}{}

\titleformat{\subsection}
  {\normalfont\large\bfseries\color{mdaDarkBlue}\sffamily}
  {\thesubsection}{0.8em}{}

\titleformat{\subsubsection}
  {\normalfont\normalsize\bfseries\color{mdaAccent}\sffamily}
  {\thesubsubsection}{0.8em}{}

% Entorno: Caja de Fuentes al inicio de cada tema
\newcommand{\cuadrofuentes}[1]{%
  \begin{tcolorbox}[
    colback=mdaLightAccent,
    colframe=mdaAccent,
    boxrule=1pt,
    arc=4pt,
    left=12pt,
    right=12pt,
    top=8pt,
    bottom=8pt,
    breakable
  ]
    \small
    \noindent\textbf{\color{mdaAccent}\large\textsf{Síntesis del Bloque y Objetivos de Evaluación:}}\\[4pt]
    #1
  \end{tcolorbox}
  \vspace{8pt}
}

% Entorno: Definición
\newtcolorbox{definicion}[1][]{
  enhanced,
  colback=mdaGray,
  colframe=mdaPrimary,
  boxrule=1pt,
  leftrule=4pt,
  arc=3pt,
  title=\textbf{\textsf{Definición Canónica\ifstrempty{#1}{}{: #1}}},
  coltitle=white,
  colbacktitle=mdaPrimary,
  attach boxed title to top left={yshift=-2mm, xshift=4mm},
  boxed title style={arc=2pt, boxrule=0pt},
  top=10pt,
  bottom=8pt,
  left=10pt,
  right=10pt,
  breakable
}

% Entorno: Ejemplo Explicado
\newtcolorbox{ejemplo}[1][]{
  enhanced,
  colback=mdaLightGreen,
  colframe=mdaGreen,
  boxrule=0.8pt,
  leftrule=3.5pt,
  arc=3pt,
  title=\textbf{\textsf{Ejemplo de Clase / Dominio\ifstrempty{#1}{}{: #1}}},
  coltitle=white,
  colbacktitle=mdaGreen,
  attach boxed title to top left={yshift=-2mm, xshift=4mm},
  boxed title style={arc=2pt, boxrule=0pt},
  top=10pt,
  bottom=8pt,
  left=10pt,
  right=10pt,
  breakable
}

% Entorno: Aviso / Ojo Examen
\newtcolorbox{aviso}[1][]{
  enhanced,
  colback=mdaLightOrange,
  colframe=mdaOrange,
  boxrule=0.8pt,
  leftrule=3.5pt,
  arc=3pt,
  title=\textbf{\textsf{Importante / Pregunta de Examen\ifstrempty{#1}{}{: #1}}},
  coltitle=white,
  colbacktitle=mdaOrange,
  attach boxed title to top left={yshift=-2mm, xshift=4mm},
  boxed title style={arc=2pt, boxrule=0pt},
  top=10pt,
  bottom=8pt,
  left=10pt,
  right=10pt,
  breakable
}

% Entorno: Tarjeta de Historia de Usuario
\newtcolorbox{tarjetahu}[1][]{
  enhanced,
  colback=mdaYellowCard,
  colframe=mdaYellowBorder,
  boxrule=1pt,
  leftrule=4pt,
  arc=3pt,
  title=\textbf{\textsf{\color{mdaDarkBlue}Tarjeta de Historia de Usuario\ifstrempty{#1}{}{: #1}}},
  coltitle=mdaDarkBlue,
  colbacktitle=mdaYellowBorder!40!white,
  attach boxed title to top left={yshift=-2mm, xshift=4mm},
  boxed title style={arc=2pt, boxrule=0.8pt, colframe=mdaYellowBorder},
  top=10pt,
  bottom=8pt,
  left=10pt,
  right=10pt,
  breakable
}

% Entorno: Regla Formal / Patrón
\newtcolorbox{teorema}[1][]{
  enhanced,
  colback=mdaLightAccent,
  colframe=mdaAccent,
  boxrule=0.8pt,
  leftrule=3.5pt,
  arc=3pt,
  title=\textbf{\textsf{Principio / Patrón Metodológico\ifstrempty{#1}{}{: #1}}},
  coltitle=white,
  colbacktitle=mdaAccent,
  attach boxed title to top left={yshift=-2mm, xshift=4mm},
  boxed title style={arc=2pt, boxrule=0pt},
  top=10pt,
  bottom=8pt,
  left=10pt,
  right=10pt,
  breakable
}

% Entorno: Ejercicio Resuelto
\newtcolorbox{ejercicio}[1][]{
  enhanced,
  colback=white,
  colframe=mdaDarkBlue,
  boxrule=1pt,
  leftrule=4pt,
  arc=3pt,
  title=\textbf{\textsf{Ejercicio de Examen Resuelto\ifstrempty{#1}{}{: #1}}},
  coltitle=white,
  colbacktitle=mdaDarkBlue,
  attach boxed title to top left={yshift=-2mm, xshift=4mm},
  boxed title style={arc=2pt, boxrule=0pt},
  top=10pt,
  bottom=8pt,
  left=10pt,
  right=10pt,
  breakable
}

% Configuración de listings para Gherkin / BDD
\lstset{
  backgroundcolor=\color{codeBg},
  basicstyle=\ttfamily\small\color{black},
  keywordstyle=\bfseries\color{mdaPrimary},
  commentstyle=\itshape\color{gray},
  stringstyle=\color{mdaGreen},
  numbers=none,
  breaklines=true,
  breakatwhitespace=false,
  showstringspaces=false,
  frame=single,
  rulecolor=\color{mdaBorderGray},
  tabsize=2,
  captionpos=b,
  inputencoding=utf8,
  extendedchars=true,
  literate=
    {á}{{\'a}}1 {é}{{\'e}}1 {í}{{\'i}}1 {ó}{{\'o}}1 {ú}{{\'u}}1
    {Á}{{\'A}}1 {É}{{\'E}}1 {Í}{{\'I}}1 {Ó}{{\'O}}1 {Ú}{{\'U}}1
    {ñ}{{\~n}}1 {Ñ}{{\~N}}1 {¿}{{?`}}1 {¡}{{!`}}1
    {«}{{\guillemotleft}}1 {»}{{\guillemotright}}1
    {€}{{\euro}}1
}
